% from ega2-1-fr.tex
\oldpage[II]{141}
\begin{proposition}[7.1.9]
\label{II.7.1.9-fr}
Soient $Y$ un préschéma localement noethérien, $p:X\to Y$ un morphisme
localement de type fini, $x$ un point de $X$, $y=p(x)$, $y'\neq y$ une
spécialisation de $y$. Il existe alors un schéma local $Y'$, spectre d'un
anneau de valuation discrète, un morphisme séparé $f:Y'\to Y$ et une
$Y$-application rationnelle $g$ de $Y'$ dans $X$, tels que, si $a$ désigne le
point fermé de $Y'$ et $b$ son point générique, on ait $f(a)=y'$,
$f(b)=y$, $g(b)=x$, et que dans le diagramme commutatif
\[
\tag{7.1.9.1}
\label{II.7.1.9.1-fr}
\begin{tikzcd}[column sep=large,row sep=large]
& k(x) \arrow[dl,"\gamma"'] \\
k(b) & k(y) \arrow[u,"\pi"] \arrow[l,"\varphi"']
\end{tikzcd}
\]
(où $\pi$, $\varphi$, $\gamma$ sont les homomorphismes correspondant
respectivement à $p$, $f$ et $g$), $\gamma$ soit une bijection.
\end{proposition}

Comme dans (\hyperref[II.7.1.4-fr]{7.1.4}), on peut se ramener au cas où $Y$
est intègre de point générique $y$ (compte tenu de
(\hyperref[I.6.3.4-fr]{I, 6.3.4, (iv)})) et comme la question est locale sur
$X$ et $Y$, on peut supposer $p$ de type fini~; on est alors dans la situation
de (\hyperref[II.7.1.4-fr]{7.1.4}), avec la propriété supplémentaire que
$k(x)$ est une extension \emph{de type fini} de $k(y)$
(\hyperref[I.6.4.11-fr]{I, 6.4.11}) et que $\mathcal{O}_{y'}$ est noethérien~;
cela permet d'appliquer (\hyperref[II.7.1.7-fr]{7.1.7}) et de prendre
$Y'=\operatorname{Spec}(A')$, où $A'$ est un anneau de valuation discrète
dominant $\mathcal{O}_{y'}$ et dont $k(x)$ est le corps des fractions. On a
donc ainsi défini un diagramme commutatif
(\hyperref[II.7.1.9.1-fr]{7.1.9.1}), où $\gamma$ est une bijection, et $\pi$
et $\varphi$ correspondent aux morphismes $p$ et $f$. En outre, comme $X$ et
$Y$ sont localement noethériens (\hyperref[I.6.6.2-fr]{I, 6.6.2}) et que $Y'$
est intègre, il y a une $Y$-application rationnelle $g$ et une seule de $Y'$
dans $X$ à laquelle correspond l'isomorphisme $\gamma$
(\hyperref[I.7.1.15-fr]{I, 7.1.15}), ce qui achève la démonstration.

\subsection{Critère valuatif de séparation}
\label{subsection:II.7.2-fr}

\begin{proposition}[7.2.1]
\label{II.7.2.1-fr}
Soient $X$, $Y$ deux préschémas, $f:X\to Y$ un morphisme quasi-compact. Les
deux conditions suivantes sont équivalentes~:
\begin{enumerate}
\item[\rm a)] Le morphisme $f$ est fermé.
\item[\rm b)] Pour tout $x\in X$ et toute spécialisation $y'$ de $y=f(x)$,
distincte de $y$, il existe une spécialisation $x'$ de $x$ telle que
$f(x')=x$.
\end{enumerate}
\end{proposition}

La condition b) exprime que $f(\overline{\{x\}})=\overline{\{y\}}$ et est donc
conséquence de a). Pour montrer que b) entraîne a), considérons une partie
fermée $X'$ de l'espace sous-jacent $X$~; soit $Y'=\overline{f(X')}$ et
montrons que $Y'=f(X')$. Considérons les sous-préschémas fermés réduits de
$X$ et $Y$ respectivement ayant pour espaces sous-jacents $X'$ et $Y'$
(\hyperref[I.5.2.1-fr]{I, 5.2.1})~; il y a alors un morphisme $f':X'\to Y'$ tel
que le diagramme
\[
\begin{tikzcd}[column sep=large,row sep=large]
X' \arrow[r,"f'"] \arrow[d] & Y' \arrow[d] \\
X \arrow[r,"f"'] & Y
\end{tikzcd}
\]
soit commutatif (\hyperref[I.5.2.2-fr]{I, 5.2.2}), et comme $f$ est
quasi-compact, il en est de même de $f'$. On est donc ramené à prouver que si
$f$ est un morphisme quasi-compact et \emph{dominant}, alors la
